%!TEX root = ../main.tex
% =============================================================================
%  一、问题重述与分析           建议 2–3 页
%  评审在问：你知道题目在问什么吗？
%  本节必须出现：① 决策时序图（图 \ref{fig:timeline}）
%                ② 问题依赖图（图 \ref{fig:dependency}）+ 误差传播表（四行齐全）
%                ③ 四问的信息结构差异：决策时刻 1→4、计划依据由确定值变预报值、
%                   两类不对称惩罚（违约 0.5 倍 / 紧急购电 5 倍）
% =============================================================================

\section{问题重述与分析}

\subsection{问题背景与决策矛盾}

\begin{guideblock}
背景只写\textbf{半页}，且只写能改变假设、约束、目标函数或验证方式的内容——
"凡是不能改变模型的行业信息，一律不写进论文"（行业内容落位原则）。
开头模板：以"计划电 \(\rightarrow\) 市场电"起笔，接新能源全面入市与取消强制配储的政策主线，
落到"储能从并网配套转变为必须依靠市场价差生存的资产"。
\textbf{禁止}空话：删掉"随着我国经济快速发展""能源互联网蓬勃发展"一类句子。
引用政策必须写全称 + 文号，且必须核对过官网原文（宁可少引三条，不可引错一条）。
\end{guideblock}

随着光伏等分布式电源大规模接入，配电网由"无源"变为"有源"，微网（microgrid）成为
促进分布式能源就地消纳、缓解并网压力的重要组织形式。本文研究的非孤岛式微网（并网型微网）
既与外部电网（以下简称外网）交换功率，又配有储能设备（BESS），其经济性主要来自
\textbf{电价的时间价差}：在电价较低或分布式能源有剩余电量的时段充电、在电价较高的时段放电，
从而以更低的购电费用满足同一小区负载。

在新能源全面参与市场交易的背景下\cite{ref136}，储能的收益来源正在发生根本变化：
依据《关于深化新能源上网电价市场化改革促进新能源高质量发展的通知》
（发改价格〔2025〕136 号），新能源上网电量原则上全部进入电力市场、电价由市场形成，
新建新能源项目也不再强制配置储能。储能的定位因此由"并网配套"转变为
\textbf{必须依靠市场价差生存的资产}，其运行策略必须直面电价与出力的双重不确定性。
\begin{guideblock}
上面这一段是本题"懂行业"的证据。若能再补一句可核对的量化对比，例如把实际谷峰价比与
储能套利的盈亏平衡阈值 \(1/(\eta_c\eta_d)\approx 1.235\) 作比较，说服力更强；
但该比值必须由本队用附件 4 自己算出，并写明时段划分口径（不同时段划分会得出不同比值）。
\end{guideblock}

本题的决策矛盾可以概括为一句话：\textbf{光伏盈余集中在电价低谷的午间，而负载高峰
落在电价尖峰的傍晚，储能容量与充放功率均有限，因此必须把"何时购电、何时充放、
预留多少安全余量"作为一个整体来优化。}这一矛盾在四个问题中被逐步放松与加难：

\begin{itemize}
  \item \textbf{问题 1（确定性计划）}：电价与负载每天相同、光伏给出一天的预测数据，
        要求 \(0{:}00\) 与 \(24{:}00\) 储电量相同，制定每天 \(0{:}00\) 的计划购电策略。
        信息结构最弱、模型最简，是整个问题的骨架。
  \item \textbf{问题 2（不确定性下的日前计划）}：电价仍固定，负载与光伏实际值逐日逐时段变化；
        购电量低于负载时须按交易时刻电价的 5 倍\textbf{紧急购电}。
        计划在 \(0{:}00\) 制定、偏差在实时结算，是一类\textbf{两阶段决策}。
  \item \textbf{问题 3（多阶段滚动调整）}：每天 \(0{:}00\)、\(6{:}00\)、\(12{:}00\)、\(18{:}00\)
        可获取未来 24 小时整点光伏预报，其他时刻可\textbf{调整}购电策略，并引入两类
        \textbf{不对称惩罚}：计划高于调整的部分按电价 0.5 倍计违约电价，调整高于计划的
        部分按 1.5 倍计。核心追问是\textbf{是否需要引入其他时刻的预报}。
  \item \textbf{问题 4（波动电价下的风险决策）}：外网电价本身实时波动，需在波动电价下
        重新求解问题 2 与问题 3，问题由"确定性价格 + 随机出力"升级为
        \textbf{价格与出力的联合不确定性}。
\end{itemize}

\subsection{四问的差别在于信息结构，而不在于模型复杂度}

\begin{guideblock}
本节必须放\textbf{决策时序图}（图 \ref{fig:timeline}），并在图下写出一句话：
"四问的差别不在于模型复杂度，而在于信息结构：决策时刻从 1 个增加到 4 个，
计划依据从确定值变为预报值，且引入了违约（0.5 倍）与紧急购电（5 倍）两类不对称惩罚。"
验收标准：评审不用读正文，只看这张图就能复述四问的差别。
\end{guideblock}

四问的数学形式看似逐问变复杂，但真正决定建模方式的只有三件事：
\textbf{决策时刻有几个}、\textbf{每个时刻知道什么}、\textbf{偏差如何被惩罚}。
图~\ref{fig:timeline} 按时间轴刻画了这三件事，并把四问放在同一坐标下对照。

\begin{figure}[htbp]
  \centering
  %!TEX root = ../main.tex
% =============================================================================
%  图 1  决策时序图（正文 §1.2）
%  规格（四项必须标注，缺一不可）：
%    ① "计划依据"与"结算依据"是两套数据，用不同线型区分
%    ② 18:00 的预报覆盖到次日 18:00，用一条跨越 24:00 的箭头表示
%    ③ 紧急购电发生在实时（结算时），不在决策时刻
%    ④ 四问在图上的差别：Q1 仅 0:00 一个决策点且无惩罚；Q2 有 0:00 + 实时惩罚；
%       Q3 有 4 个决策点；Q4 旁注"电价亦为随机"
% =============================================================================
\begin{tikzpicture}[
  x=0.46cm, y=1cm, font=\small, >=Stealth,
  ev/.style={font=\scriptsize, inner sep=1.5pt},
]

% ---------- 泳道底色 ----------
\fill[gray!8, rounded corners=2pt] (-3.15,-2.15) rectangle (26.2,-0.30);   % 信息
\fill[blue!5, rounded corners=2pt] (-3.15,-3.70) rectangle (26.2,-2.30);   % 决策
\fill[red!4,  rounded corners=2pt] (-3.15,-5.25) rectangle (26.2,-3.85);   % 结算

% ---------- 泳道标签 ----------
\node[rotate=90, font=\footnotesize\heiti, gray!60!black] at (-2.55,-1.22) {信息};
\node[rotate=90, font=\footnotesize\heiti, gray!60!black] at (-2.55,-3.00) {决策};
\node[rotate=90, font=\footnotesize\heiti, gray!60!black] at (-2.55,-4.55) {结算};

% ---------- 时间刻度（横轴 0:00 -> 24:00） ----------
\foreach \h in {0,6,12,18,24}{
  \draw[gray!25] (\h,-5.25) -- (\h,-0.30);
}
\foreach \h/\lab in {0/0:00, 6/6:00, 12/12:00, 18/18:00, 24/24:00}{
  \draw[gray!60] (\h,0.00) -- (\h,-0.16);
  \node[above, font=\scriptsize, gray!70!black] at (\h,0.00) {\lab};
}
\draw[gray!60, ->] (-0.9,0.00) -- (26.2,0.00) node[right, font=\scriptsize, gray!70!black] {时刻};
\draw[gray!45, dashed] (24,0.00) -- (24,-5.25);
\node[font=\scriptsize, gray!70!black, align=center] at (25.45,-5.55) {次日};

% ---------- 泳道 1：信息（附件 3 的四个预报发布时刻） ----------
% 预报覆盖区间的起点用实心圆标注
\foreach \h/\y/\len/\txt in {%
  0/-0.62/24/{\,覆盖至 24:00},
  6/-0.95/18/{\,覆盖至 24:00},
  12/-1.28/12/{\,覆盖至 24:00},
  18/-1.61/8/{~覆盖至次日 18:00（跨日）}}{
  \fill[blue!55!black] (\h,\y) circle (1.5pt);
  \draw[gray!60, line width=0.5pt, -{Stealth[length=4pt]}] (\h,\y) -- (\h+\len,\y);
  \node[ev, anchor=west, gray!25!black] at (\h+0.35,\y) {\h:00 发布\txt};
}

% ---------- 泳道 2：决策 ----------
% Q1：唯一决策时刻 0:00
\fill[blue!65!black] (0,-2.62) circle (2.4pt);
\node[font=\scriptsize, anchor=west, blue!50!black] at (0.35,-2.62)
  {计划购电：Q1\,/\,Q2\,/\,Q3\,/\,Q4（唯一决策时刻，随后不再修正）};
% Q3/Q4：6:00、12:00、18:00 调整
\foreach \h in {6,12,18}{
  \draw[blue!65!black, fill=white, line width=0.7pt] (\h,-3.28) circle (2.4pt);
}
\draw[gray!55, dotted, line width=0.6pt] (0,-2.62) -- (0,-3.28);
\node[font=\scriptsize, anchor=west, blue!50!black] at (6.35,-3.28)
  {调整购电：Q3\,/\,Q4 在 6:00、12:00、18:00 各修正一次};

% ---------- 泳道 3：结算 ----------
\draw[gray!55, densely dashed, line width=0.7pt, -{Stealth[length=5pt]}]
  (0,-4.20) -- (24,-4.20);
\node[font=\scriptsize, anchor=west, gray!25!black] at (0.35,-3.98)
  {结算依据：附件 2 实际负载与实际光伏功率（\textbf{与计划依据是两套数据}）};
% 紧急购电：发生在实时（结算时），不在决策时刻
\foreach \x in {9.4, 15.1, 21.3}{
  \draw[red!75!black, line width=0.9pt] (\x-0.28,-4.52) -- (\x+0.28,-4.52) -- (\x,-4.78) -- cycle;
}
\node[font=\scriptsize, anchor=west, red!65!black] at (-3.05,-4.98)
  {缺口 $\rightarrow$ 紧急购电（按 $5p_t$ 结算），发生在\textbf{实时}，不在任何决策时刻};

% ---------- 图例 ----------
\node[ev, anchor=west] at (-3.05,-5.62) {%
  \tikz\draw[gray!60,line width=0.5pt,-{Stealth[length=4pt]}] (0,0)--(0.55,0);
  ~预报（计划依据）\quad
  \tikz\draw[gray!55,densely dashed,line width=0.7pt,-{Stealth[length=4pt]}] (0,0)--(0.55,0);
  ~实际值（结算依据）\quad
  \tikz\fill[blue!65!black] (0,0) circle (2.2pt);~决策时刻\quad
  \tikz\draw[red!75!black,line width=0.9pt] (0,-0.05)--(0.28,-0.05)--(0.14,-0.26)--cycle;
  ~紧急购电（非决策时刻）\quad
  注：Q4 中电价亦为随机（旁注）};
\end{tikzpicture}

  \caption{微网决策时序与信息结构（横轴为一天 24 小时；上行泳道为信息到达时刻，
           中行为决策时刻，下行为结算依据；\(\Delta t=1/6\)~h，共 144 个时段）}
  \label{fig:timeline}
\end{figure}

\begin{guideblock}
图 \ref{fig:timeline} 的四个"必须标注"（缺一不可）：
① "计划依据"与"结算依据"是\textbf{两套数据}，用不同线型区分；
② \(18{:}00\) 的预报覆盖到次日 \(18{:}00\)，用一条跨越 \(24{:}00\) 的箭头表示；
③ 紧急购电发生在\textbf{实时（结算时）}，不在任何决策时刻；
④ 四问在图上的差别：Q1 仅 \(0{:}00\) 一个决策点且无惩罚；Q2 有 \(0{:}00\) + 实时惩罚；
   Q3 有 4 个决策点；Q4 旁注"电价亦为随机"。
\end{guideblock}

由图~\ref{fig:timeline} 可直接读出四问在信息结构上的递进：问题 1 只有一个决策时刻且
"知道当日真值"；问题 2 把决策时刻固定在 \(0{:}00\)、把已知量换成预报值，并在实时结算时
用 5 倍电价惩罚缺口；问题 3 把决策时刻扩展到 4 个，但计划量基准始终是 \(0{:}00\) 的计划值
（假设 A7），形成带\textbf{非预期性约束}的多阶段滚动结构；问题 4 进一步让价格本身随机。

\subsection{四问的依赖关系与误差传播}

\begin{guideblock}
本节必须放\textbf{问题依赖图}（图 \ref{fig:dependency}）与\textbf{误差传播表}
（表 \ref{tab:errorprop}，四行齐全且"本文的处理"列非空）。目的：证明四问是一条链，
而不是四篇互不相干的报告。
\end{guideblock}

\begin{figure}[htbp]
  \centering
  %!TEX root = ../main.tex
% =============================================================================
%  图 2  四问的模型依赖关系（正文 §1.3）
%  规格：附件 -> Q1 -> Q2 -> Q3 -> Q4，每个箭头必须标注"提供什么"；
%        本图用于证明四问是一条信息结构递进链，而不是四篇互不相干的报告。
% =============================================================================
\begin{tikzpicture}[
  font=\small, >=Stealth,
  att/.style ={draw=gray!50, fill=gray!8, rounded corners=2pt, align=center,
               font=\scriptsize, text width=2.0cm, minimum height=0.72cm, inner sep=3pt},
  pre/.style ={draw=orange!55!black, fill=orange!8, rounded corners=2pt, align=center,
               font=\scriptsize, text width=6.6cm, minimum height=0.85cm, inner sep=4pt},
  qbox/.style={draw=blue!50, fill=blue!7, rounded corners=2pt, align=center,
               font=\small, text width=4.6cm, minimum height=0.95cm, inner sep=4pt},
  obox/.style={draw=green!45!black, fill=green!7, rounded corners=2pt, align=center,
               font=\scriptsize, text width=2.2cm, minimum height=0.8cm, inner sep=3pt},
  ar/.style  ={->, draw=gray!65, line width=0.6pt},
  lab/.style ={font=\scriptsize, anchor=west, align=left, text width=5.6cm,
               inner sep=1pt, gray!20!black},
]

% ---------------- 附件 ----------------
\node[att] (a1) at (-4.35,0) {附件 1\\电价/负载/光伏};
\node[att] (a2) at (-1.45,0) {附件 2\\负载/光伏实际值};
\node[att] (a3) at ( 1.45,0) {附件 3\\四个时刻光伏预报};
\node[att] (a4) at ( 4.35,0) {附件 4\\波动电价};

% ---------------- 数据预处理 ----------------
\node[pre] (pre) at (0,-1.35)
  {数据预处理：时标对齐（标签 $=$ 时段末端）\ \textbullet\ 附件 1 身份核查\ \textbullet\ 特征与误差分析（第三章）};

\foreach \n in {a1,a2,a3,a4}{\draw[ar] (\n.south) -- (pre.north);}

% ---------------- 四问 ----------------
\node[qbox] (q1) at (0,-2.95) {\textbf{问题 1}\ 确定性线性规划（LP）};
\node[qbox] (q2) at (0,-4.60) {\textbf{问题 2}\ 两阶段随机规划};
\node[qbox] (q3) at (0,-6.25) {\textbf{问题 3}\ 多阶段滚动 $+$ 非预期性};
\node[qbox] (q4) at (0,-7.90) {\textbf{问题 4}\ 波动电价 $+$ 风险度量};

\draw[ar] (pre.south) -- (q1.north);

% 箭头上的"提供什么"标注
\draw[ar] (q1.south) -- (q2.north);
\node[lab] at (0.35,-3.78)
  {\textbf{提供}\ 功率平衡与储电量递推骨架；无储能基线 $C^{base}$ 口径；
   储电水价值与阈值律语言};

\draw[ar] (q2.south) -- (q3.north);
\node[lab] at (0.35,-5.43)
  {\textbf{提供}\ 固定场景集（消融与 VoI 共用）；报童临界分位数给出的安全边际；
   紧急购电的结构性认识};

\draw[ar] (q3.south) -- (q4.north);
\node[lab] at (0.35,-7.08)
  {\textbf{提供}\ 多阶段滚动框架与非预期性实现；终端条件处理；
   预报频次对照（VoI）方法};

% ---------------- 输出 ----------------
\node[obox] (o1) at (-5.55,-2.95) {result1.xlsx\\+ 表 1 / 表 2};
\node[obox] (o2) at (-5.55,-4.60) {result2.xlsx\\+ 表 3};
\node[obox] (o3) at (-5.55,-6.25) {result3.xlsx\\+ VoI 结论};
\node[obox] (o4) at (-5.55,-7.90) {result4-2.xlsx\\result4-3.xlsx};

\foreach \q/\o in {q1/o1,q2/o2,q3/o3,q4/o4}{
  \draw[ar, draw=green!45!black] (\q.west) -- (\o.east);
}

\end{tikzpicture}

  \caption{四问的模型依赖关系与复用内容（箭头标注上游向下游提供的模型要素）}
  \label{fig:dependency}
\end{figure}

图~\ref{fig:dependency} 表明四问共享同一套物理模型骨架（功率平衡、储电量递推、上下界约束），
区别仅在于信息结构与风险处理方式：问题 1 提供模型骨架、水价值语言与无储能基线 \(C^{base}\)；
问题 2 加入场景与两阶段结构，产出安全边际与紧急购电结构；问题 3 复用问题 2 的场景集以保证
可比性，产出滚动框架与信息价值结论；问题 4 复用问题 3 的多阶段结构而把电价改成随机过程。
这种"逐层加模块"的组织方式使每一层的新增贡献都可被\textbf{单独消融}（见第九章）。

\begin{table}[htbp]
  \centering
  \caption{四问之间的误差传播路径与本文的处理方式}
  \label{tab:errorprop}
  \footnotesize
  \begin{tabularx}{\linewidth}{L{2.0cm} L{3.4cm} L{4.2cm} Y}
    \toprule
    上游 \(\rightarrow\) 下游 & 传递的内容 & 潜在的误差传播 & 本文的处理 \\
    \midrule
    附件 \(\rightarrow\) 问题 1
      & 电价、负载、光伏预测的时标与单位
      & 时标错位一个槽位或丢失 \(0{:}00{+}1\) 列，会使全部合计系统性偏离
      & 统一按"标签 = 时段末端"对齐，并断言每天恰有 144 个时段（§3.2） \\
    \addlinespace
    问题 1 \(\rightarrow\) 问题 2
      & 模型骨架、\(C^{base}\) 口径、水价值语言
      & 问题 1 用"全年平均日"、问题 2 用"334 天窗口"，两套口径若混用，节省率不可比
      & 显式声明两套口径并各自命名；四问统一用同一 \(C^{base}\) 定义（\calref{C-01}） \\
    \addlinespace
    问题 2 \(\rightarrow\) 问题 3
      & 场景生成与削减方法、安全边际
      & 场景分布偏离真实分布，若各问各生成场景，"改进"可能只是场景噪声
      & 场景集\textbf{一次性固定并存盘}，问题 3 与消融实验读同一份场景集（§9.4） \\
    \addlinespace
    问题 3 \(\rightarrow\) 问题 4
      & 多阶段滚动框架、终端条件
      & 电价随机性与预报误差叠加，若分别抽样会丢失价格—出力的相关结构
      & 用分块 Bootstrap 生成\textbf{联合}情景，保留 \(\rho(p,L)\) 相关结构（§7.2） \\
    \bottomrule
  \end{tabularx}
\end{table}

\begin{guideblock}
误差传播表的验收：四行齐全，"本文的处理"列非空，且每条都能指向正文中的具体小节。
最后一行的关键点：\textbf{四问必须统一使用同一基线定义 \(C^{base}\)}，
否则论文中出现的多个"节省率"互不可比，这是最容易被评审抓住的内部矛盾。
\end{guideblock}
